\section{A gate of bounded cost for every alphabet}\label{sec:recurrence}

The alphabet enters the mask recursion only through the number \(b\)
of bits needed to name a letter. We will obtain three successive cost
bounds in \(b\). On the inner anchors, of surplus \(2\rho\), a linear
bound with slope \(\rho/4\) suffices. A call on these anchors replaces a
mask diagonal and saves \(2\rho\) per remaining bit before paying for the
call. This saving gives a logarithmic bound on the original anchors, of
surplus \(\rho\). Finally, the small increase in each group's width
given by Lemma~\ref{lem:layout} makes the cost bounded independently of
\(b\).

\subsection{Two estimates for the recursive test}

Use the overheads \(g_0,G\) of Lemma~\ref{lem:mask-recursion}. The
unchanged active set costs \(g_0\) at a recursive step. A call on a
smaller active set must also pay \(G\) to install its split data. For
\(n\) active satellites put
\[
 D(n)=\left\lfloor\frac{\log n}{-\log(1-p_0)}\right\rfloor,\qquad
 l_0=\lceil4g_0/\rho\rceil,\qquad
 A_j=10(l_0+G)(j+1).
\]
Every eligible smaller set has \(D(n')\le D(n)-1\), since
\(n'\le(1-p_0)n\). Put \(D_*=\lceil H\rceil+D(r)\) and \(A_*=A_{D_*}\).
In the binary layout \(p_0=1/256\), so \(D_*=O(H)\). The additional
\(\lceil H\rceil\) in this definition ensures
\(A_*\ge (g_0/\rho)H\).

\begin{lemma}[Linear and logarithmic costs]\label{lem:two-rates}
For every \(b\ge0\) and \(2\le n\le r-2\), the gates of
Section~\ref{sec:masks} on the groups of surplus \(2\rho\) and \(\rho\)
admit the respective net cost bounds
\begin{align}
 f_{2\rho}(b,n)&\le A_{D(n)}+\rho b/4, \label{eq:inner-rate}\\
 f_\rho(b,n)&\ll A_*+(g_0/\rho)\log(2+\rho b/A_*).
 \label{eq:root-rate}
\end{align}
\end{lemma}
\begin{proof}
For \eqref{eq:inner-rate}, induct on \(b\), simultaneously for all
\(n\), and set \(D=D(n)\). The exact comparison of
Lemma~\ref{lem:mask-recursion} costs \(2\rho b\), which proves the
claim when \(7\rho b/4\le A_D\). Above this threshold we have
\(l_0\le b\). Take \(l=l_0\), \(b'=b-l\) in \eqref{eq:recursion}.
In the main branch the linear bound decreases by
\(\rho l/4\ge g_0\), enough to pay the overhead.

If a correction has an eligible child, \(D\ge1\), and the cost
increment in \eqref{eq:recursion} for each head using such a child
is at most
\[
  3l+g_0+G+A_{D-1}-7\rho b'/4
  < (3+7\rho/4)l+g_0+G-(A_D-A_{D-1})<0 .
\]
Here \(A_D-A_{D-1}=10(l+G)\). At least one head must use the child.
The largest upper bound therefore occurs with one. Including the
outside costs \(g_0+2\rho b'\), the correction costs at most
\[
 A_{D-1}+3l+2g_0+G+\rho b'/4\ \le\ A_D+\rho b/4 .
\]
This proves the first estimate, including the case with no eligible child.

For \eqref{eq:root-rate}, use the exact comparison while
\(\rho b\le100A_*\). At greater lengths take
\(l=\lceil\rho b/50\rceil\) and \(b'=b-l\). By \eqref{eq:inner-rate}
the increment for each head using a child is at most
\[
  3l+g_0+G+A_*-7\rho b'/4\ \le\ -\rho b .
\]
Indeed \(l\le\rho b/50+1\), \(g_0+G\le A_*\), and
\(0<\rho<10^{-3}\). This pays even the full outside charge \(\rho b'\)
of the correction. Equation~\eqref{eq:recursion} thus gives
\(f_\rho(b,n)\le g_0+f_\rho(b',n)\), using nonnegative cost bounds.
Each such step leaves \(b'\le(1-\rho/50)b\) bits. It takes at most
\(O(\rho^{-1}\log(2+\rho b/A_*))\) steps to reach the exact comparison,
whose cost there is at most \(100A_*\).
\end{proof}

\subsection{Removing the dependence on the bit length}

Enlarge each of the \(r\) groups in the layout of surplus \(\rho\) by
the width increment \(\alpha=\rho/(10r)\) of Lemma~\ref{lem:layout}.
Every satellite participates in this final construction. We can still
call the gate on the original groups, with cost bound
\(f(t)=f_\rho(t,r-2)\) from \eqref{eq:root-rate}. Its activity flags
are all one, so its split data are fixed. The wider groups give the
exact checks a saving proportional to \(t\), which will dominate
\(f(t)\).

\begin{lemma}[Uniform cost in the alphabet]\label{lem:uniform-mask}
On the enlarged groups, for every \(b\ge0\) there is a gate of net
cost \(F(b)\ll A_*\) satisfying Lemma~\ref{lem:diagonal} on the full
domain of replicated \(b\)-bit letters. It returns a fixed nonzero
tensor on a constant column and zero on every column with a repeated
letter and at least two different letters.
\end{lemma}
\begin{proof}
At \(b=0\) the unit gate has \(F(0)=0\). For the step from \(t\)
to \(t+1\) bits, write each letter as \((a,v)\), where
\(a\in\{0,1\}\), \(v\in\mathbb F_2^t\). Supply two independent
full mask diagonals with labels \(m_0,m_1\in\mathbb F_2^t\).
Keep both mask slots and apply the shorter enlarged-group gate to the
letters \(v+m_a\). On a constant column its return is both full
diagonals times the shorter gate's offset \(Q\). Its net cost against
this offset is \(F(t)\).

On an inadmissible column a repeated original letter stays repeated
after masking. Thus a mask pair contributes only when every masked
tail equals one common value \(z\). Both heads must then occur.
Otherwise adding the same mask to every tail preserves the original
disagreement. We reproduce these surviving terms exactly, without the
guard of Section~\ref{sec:guard}, so the repeat may involve an anchor.

\paragraph{The two occupied heads.}
Sum over the two anchor head bits and check them locally. If they
are equal, also sum over the first satellite of the other head, checking
its head and checking the opposite bit at every earlier satellite.
When both heads occur exactly one branch passes, and when only one
occurs none passes. The logarithm of the number of branches is \(O(H)\).
In a passing branch sum over a public \(t\)-bit label \(z\). For each
\(a\in\{0,1\}\), the exact mask output \(I_a(z)\) is \(E_S(v+z)\)
if all letters with head \(a\) have a common tail \(v\), and is zero
otherwise. Since both heads are occupied, the tensor
\(\sum_z I_0(z)\otimes I_1(z)\) supplies exactly the mask pairs that
produce a common masked tail. Every such pair occurs once.

Let \(u_a\) be the sum of the widths of the \emph{original} anchors
with head \(a\), so \(u_0+u_1=1-\rho\). The exact output \(I_a(z)\)
can be produced at cost at most \((1-u_a-\alpha)t\).
With one anchor of head \(a\), broadcast its tail from the enlarged
anchor. With no anchor of this head, broadcast from its designated
enlarged satellite, of width \(\alpha\). In either case compare the
broadcast tail with every member of the head and add \(z\) locally
to form the mask. If both anchors have head \(a\), first compare their
tails at cost \((\rho-2\alpha)t\). On the surviving domain their
union, of width one, broadcasts the common tail at zero cost.
These operations use the domain version of Lemma~\ref{lem:broadcast}.

\paragraph{An alternative output at a repeated head.}
Construct \(T_a(z)\) on the original groups. Give each member of head
\(a\) its letter \(v+z\), and independently sum over a dummy \(t\)-bit
letter for each member of the other head. A dummy satellite is summed
at its single original site by the local vector \(\sum_v e_v\).
A dummy anchor requires a partial diagonal, so the dummy cost is
\((1-\rho-u_a)t\). Each satellite, real or dummy, takes part in the
call of cost \(f(t)\). Its original site can choose locally between
the real letter and the dummy sum. Thus the call uses the fixed split
data for all \(r-2\) satellites.

Apply the promise broadcast from the two original anchors, keep its
output in the mask slot, and erase the temporary registers. When the
anchors agree, every site receives their common letter.
Lemma~\ref{lem:broadcast}(iii) bounds the net cost on the whole
domain by \(O(\log(2|S|))\). The gate and broadcast offsets have a fixed
product \(P_a\), padded across public branches if necessary.
Thus \(T_a(z)\) has net cost at most
\[
       (1-\rho-u_a)t+f(t)+O(\log(2|S|)).
\]
Pad the exact output to \(\widehat I_a(z)=I_a(z)\otimes P_a\).

Suppose a pair of equal original letters has head \(a\). The input
of the original-group gate then already contains a repeat, whatever
the dummies are. It can survive only if all its inputs agree.
If two real tails in this head differ there is no surviving term.
If every real tail has value \(v\), the unique surviving assignment
sets every dummy to \(v+z\), and the broadcast acts on its promise.
In either case \(T_a(z)=\widehat I_a(z)\).

For each \(z\) use the correction
\[
 \widehat I_0(z)\otimes\widehat I_1(z)
 -
 \big(\widehat I_0(z)-T_0(z)\big)\otimes
 \big(\widehat I_1(z)-T_1(z)\big).
\]
At every inadmissible column with both heads present the second
product is zero, since one head contains a repeat. This conclusion
holds for every output of the other head's broadcast away from its
promise. Summing the first product over \(z\) therefore recovers
the surviving terms of the main branch. Append \(Q\) to the
correction and \(P_0\otimes P_1\) to the main branch, then subtract.
On inadmissible columns with a single head both branches already
vanish. On a constant column the correction vanishes and the return
is \(G_S(2^t)^{\otimes2}\otimes Q\otimes P_0\otimes P_1\).

\paragraph{Summing the costs.}
The product of the two exact outputs cancels in the expansion of the
correction. Every remaining summand uses one or two of the cheaper
outputs \(T_a\). Relative to the price \(t+h(P_a)\) of its slot in
the common offset, an exact head saves at least \((u_a+\alpha)t\).
A head using \(T_a\) contributes at most
\(f(t)-(u_a+\rho)t+O(\log(2|S|))\). Pay \(t\) for the shared \(z\).
Since \(u_0+u_1=1-\rho\), the resulting recurrence is
\begin{equation}\label{eq:outer-sum}
 F(t+1)\le
 \operatorname{lsum}\!\left(F(t),\;
 Cg_0+\max\{f(t)-\alpha t,\;2f(t)-\rho t\}\right),
\end{equation}
where \(\operatorname{lsum}(x,y)=\log_2(2^x+2^y)\) and \(C\) is
absolute. The first correction cost uses one \(T_a\), the second
uses both. All public choices and the sum of the three surviving
terms are included in \(Cg_0\).

Set \(x=\rho t/A_*\). Since \(A_*\ge(g_0/\rho)H\) and
\(\log r\le H\), the logarithmic term in \eqref{eq:root-rate} costs
\(O(A_*)\) up to \(x=r\). For all \(x\ge0\), use
\(\log(2+x)\le\log r+\log(2+x/r)\) and
\(\log(2+y)\le C_\delta+\delta y\). Choosing \(\delta>0\) small
enough gives
\[
                     f(t)\le C_1 A_*+\alpha t/8 .
\]
Both correction costs in \eqref{eq:outer-sum} are therefore at most
\(C_2 A_*-\alpha t/2\), including the overhead. Iterating after
exponentiation gives
\[
  2^{F(b)}
     \le 1+2^{C_2 A_*}\sum_{t=0}^{b-1}2^{-\alpha t/2}
     \le 1+O(2^{C_2 A_*}/\alpha).
\]
Finally \(\log(1/\alpha)=\log(10r/\rho)=O(A_*)\) by
Lemma~\ref{lem:layout}. This proves the bound uniformly in \(b\).
\end{proof}

\subsection{The bound for sets}

\begin{proof}[Proof of Theorem~\ref{thm:main}]
Let \(r\ge4\), and take a nonempty sunflower-free family of \(N\)
words of length \(k\). Embed the alphabet of each coordinate in a
binary cube, choosing the length independently at each coordinate.
Lemma~\ref{lem:uniform-mask} gives a gate of net cost \(O(A_*)\)
for each cube. The enlarged groups have surplus \(9\rho/10\).
Lemma~\ref{lem:diagonal} now gives
\[
                         (9\rho/10)\log_2 N\ll k A_* .
\]
In the binary layout \(d\ll H\), \(\ell\ll L\), \(\log |S|\ll H\),
and \(\rho^{-1}\ll H\sqrt L\) by Lemma~\ref{lem:layout}. The split
cost \(G=O(dqH)\) of Lemma~\ref{lem:splitdata} is \(O(H^2)\), and
the guard and summation overheads give
\[
 g_0/\rho\ll H/\rho+\ell/\rho^2\ll H^2 L^2,\qquad
 A_*\ll H(g_0/\rho+G)\ll H^3L^2 .
\]
Hence \(\log_2 N\ll kH^4L^{5/2}\). For \(r=3\) the fixed planar
construction in Section~\ref{sec:plane} gives \(\log_2 N=O(k)\).

For a \(k\)-uniform set family the coloring in Section~\ref{sec:words}
retains at least an \(e^{-k}\) fraction of its members as words.
Multiplication by \(e^k\) is absorbed by the absolute constant in
\(B_r\). For sets of size at most \(k\), pad each member to size \(k\)
with fresh elements belonging only to that member. Padding preserves
all intersections between distinct members, so the bound applies to
the original family as well.
\end{proof}

The diagonal comparison also applies when a word family contains
sunflowers, provided a tensor on the word indices can annihilate their
contributions. This costs only once for a whole tuple of words.

\begin{corollary}[A tensor vanishing on sunflowers]\label{cor:weighted}
Let \(r\ge3\) and let \(\mathcal W\) consist of \(N\ge1\) distinct words
of length \(k\). Suppose an order-\(r\) array over the working field,
indexed by \(\mathcal W^r\), has tensor rank \(R\), equals one on the
diagonal, and vanishes on every nondiagonal tuple admissible in
every coordinate. Then
\[
        \log_2 N\ll kH^4L^{5/2}+H\sqrt L\,\log_2 R .
\]
\end{corollary}
\begin{proof}
Multiply each joint word-index term in the diagonal comparison by
the specified array entry. For a simple tensor, one representative
of each group multiplies by its factor locally. Summing over a rank
decomposition costs at most \(\log_2 R\). The coordinate gates kill
every tuple with an inadmissible column. The array kills the remaining
nondiagonal tuples, and its diagonal entries preserve the common
offset. Thus \((9\rho/10)\log_2 N\ll k A_*+\log_2 R\) for \(r\ge4\).
For \(r=3\) apply the same weighting to the planar comparison of
Proposition~\ref{prop:comparison}, with the three groups of
Section~\ref{sec:plane}. Their positive surplus and the total
cost per coordinate are absolute constants, giving
\(\log_2 N\ll k+\log_2 R\).
\end{proof}

The optimal base \(b_r=\sup_{k\ge1}M_r(k)^{1/k}\) is also the limiting
growth rate as \(k\to\infty\) by the product construction of
Abbott~\cite[Theorem~1]{Abbott66}. Take sunflower-free uniform
families of respective set sizes \(k,l\) on disjoint ground sets, and
form all unions of a member of each. If \(r\) of the unions formed a
sunflower, their projections to either ground set would have constant
pairwise intersection. If two projections there were equal, that
common intersection would be the entire projected set. By uniformity
every projection containing it would equal it. If no two were equal,
the projections would themselves be a sunflower. Thus both projections
would be constant, contradicting distinctness of the unions.
Consequently \(M_r(k+l)\ge M_r(k)M_r(l)\), and Fekete's
lemma applied to \(\log M_r(k)\) gives
\[
                         b_r=\lim_{k\to\infty}M_r(k)^{1/k}.
\]
The block transversals in the introduction give \(b_r\ge r-1\).
Whether \(b_r=O(r)\) remains open.
